% =====================================================================
%  2. Área de abrangência e critérios de seleção
% =====================================================================
\section{Área de abrangência e critérios de seleção}
\label{sec:area}

A ação abrange todo o território do Estado do Paraná, contemplando as
Microrregiões de Água e Esgoto. Isso não significa que todos os municípios
serão atendidos. Os territórios e as famílias beneficiadas serão escolhidos
pela aplicação de critérios definidos em âmbito estadual, que variam conforme
a frente de trabalho \cite{mop2026}.

Para as comunidades que receberão sistemas coletivos de abastecimento de água
e, na mesma área, sistemas descentralizados de esgotamento sanitário, o
primeiro critério é o adensamento populacional. Serão consideradas as
comunidades, ou os conjuntos de comunidades, com 70 ou mais domicílios
ocupados, número a partir do qual se espera viabilidade técnica e econômica
compatível com os modelos de gestão a serem estabelecidos. O segundo critério
é o nível de desenvolvimento do município, medido pelo Índice Ipardes de
Desempenho Municipal -- IPDM \cite{ipdm}. O terceiro é a existência de
associações de moradores já constituídas para a gestão de sistemas de
abastecimento de água que necessitem de complementação ou melhorias, o que
favorece a reabilitação de sistemas em funcionamento e aproveita a
organização comunitária existente.

Para as áreas que receberão apenas sistemas descentralizados de esgotamento
sanitário, os critérios estão ligados à proteção dos mananciais. Serão
priorizados os mananciais prioritários do PSH-PR e os mananciais de
abastecimento público selecionados por critérios quali-quantitativos em que o
IDR-Paraná atuará, considerando os domicílios situados em área rural segundo o
Censo Demográfico 2022 \cite{ibge2022}.

A \autoref{fig:selecao} mostra como os dois conjuntos de critérios levam às
duas frentes de implantação descritas nas \autoref{sec:abastecimento} e
\ref{sec:esgotamento}.

\begin{figure}[htbp]
  \centering
  \caption{Aplicação dos critérios de seleção e definição das frentes de trabalho}
  \label{fig:selecao}
  \fluxograma{fluxograma-selecao}
  \fonte{Elaborado pelo IAT (2026) com base em \citeonline{mop2026}.}
\end{figure}

\pendencia{Os critérios não indicam pesos, limiares (exceto os 70
domicílios) nem a ordem de aplicação. Para a reprodutibilidade da seleção,
convém descrever o procedimento de ponderação ou hierarquização e o
responsável por aplicá-lo. O POA 2026--2027 prevê como produto dessa etapa
um mapa dos municípios selecionados (Anexo~\ref{an:cronograma}).}
